\section{模型检验与敏感性分析}

为验证模型的有效性与结果的可靠性，本章从\textbf{独立复算校验}、\textbf{最优性论证}、\textbf{参数敏感性}与\textbf{误差来源}四个方面进行检验。

\subsection{独立可行性校验器}\label{sec:verify}

本文建立了\textbf{独立于求解器}的可行性校验机制，其设计原则是：校验器\textbf{不导入求解器的任何模块}，而是从附件参数与 DEM 出发\textbf{从零重算}全部物理量（航段时间、能耗、SOC、载荷、资源占用区间），再与方案上报值逐条比对。校验项覆盖载质量、装载体积、能量预算、返航 SOC、时限、资源数量、资源时段、通信覆盖等 \textbf{18 类}约束。

\textbf{负样本测试是硬要求}：对故意构造的违规方案（删减货箱、篡改上报能耗与 SOC、超载、资源时段重叠、通信中断未覆盖等），校验器\textbf{必须报错}。只测通过用例的校验器等于没有校验器。

\textbf{校验结果}：问题二\textbf{硬违规 0 条}；问题三\textbf{硬违规 0 条}，另有 42 条\textbf{软违规}，全部属于“中继资源不足”这一\textbf{题目内在缺口}（见 \S\ref{sec:q3gap}）。本文特意把两者分为两类：\textbf{资源缺口}如实上报、不阻断产出；\textbf{排班缺陷}（已安排中继却未盖住中断区间）仍作为硬违规拦截。

\subsection{最优性论证}

\textbf{（1）问题一可证最优。}通过 Martello--Toth 型解析下界（式~\eqref{eq:lb}）证明 18 架次\textbf{在架次数维度上最优}——15 个服务区的启发式架次数\textbf{全部等于其下界}。这是“下界 $+$ 可行构造”构成最优性证明的标准形式。

\textbf{（2）问题二未证最优。}25 架次距理论下界 24 仅高 1 个，但本文只称其为“经验证的可行解”。\textbf{本文不把启发式结果称为全局最优}。

\textbf{（3）问题四为精确结论。}合法组数判据 $K\in[1,3]$ 是\textbf{精确的图论计算}结果，不含启发式近似。

\subsection{敏感性分析}

\begin{table}[htbp]
	\caption{敏感性分析汇总（本文结果）}
	\label{tab:sens}
	\centering
	\zihao{-5}
	\begin{tabular}{p{3.1cm}p{6.0cm}p{5.2cm}}
		\toprule
		参数 & 敏感性结论 & 依据 \\
		\midrule
		返航安全余量 $\rhog$
		& \textbf{高敏感}。$\rhog$ 由 $0.20$ 增至 $0.35$ 使架次 $18\to25$、能耗由 \SI{75.07}{kWh} 升至 \SI{106.60}{kWh}；$\rhog\ge0.375$ 时\textbf{无可行解}
		& 表~\ref{tab:q1rho}；远距离服务区空载亦无法返航 \\
		通信采样步长 $\Delta t$
		& \textbf{中等敏感}。\SI{5}{s} 细化到 \SI{2}{s} 时中断占比判定变化小于 2\%；\SI{30}{s} 会漏判约 8\% 短时中断
		& 遮挡时段多为秒级，步长过粗会漏判 \\
		悬停网格步长
		& \textbf{中等敏感}。$200\sim800$\,m 时\textbf{几何可达}覆盖率均达 100\%；\SI{1500}{m} 降至约 86\%
		& 仅影响选址能力上界，不改变资源缺口 \\
		衰落裕量 $M$ / DEM 高程噪声
		& \textbf{不敏感}。$M$ 由 \SI{2}{dB} 增到 \SI{18}{dB} 在合理范围内覆盖率不变；$\sigma=\SI{10}{m}$ 高程扰动下能耗偏移 $<1\%$
		& 选址留有较大裕度；爬升能耗占比有限 \\
		\bottomrule
	\end{tabular}
\end{table}

\subsection{误差来源}

\textbf{（1）DEM 为 DSM 带来的保守偏差。}附件 DEM 为数字表面模型（含植被与建筑），高程 $\ge$ 真实地形，故净空与遮挡判定\textbf{偏保守}——这对救援场景的安全性有利，但可能低估通信可用性。

\textbf{（2）单位换算风险。}自由空间损耗公式中 $f$ 为 MHz、$D$ 为 km，而本文内部统一使用 SI 单位（距离以 m 存储）。代入前必须除以 1000，否则产生约 \SI{60}{dB} 的系统性偏差。本文在链路层集中处理该换算并加单元测试锁定。

\textbf{（3）数值反解误差。}最大安全载荷由 Brent 法反解，实测回代相对误差 $5.6\times10^{-15}$，可忽略。

\textbf{（4）时间离散误差。}通信判定按 $\Delta t=\SI{2}{s}$ 采样，理论上会漏掉短于一个步长的中断；敏感性分析表明该步长足以捕捉遮挡特征。

\subsection{两次被纠正的建模缺陷及其传导}\label{sec:changes}

本文在求解过程中发现并纠正了两处会\textbf{改变结论方向}的建模缺陷。由于四问是链式继承关系，每次修正都向\textbf{下游全部传导重算}：

\begin{table}[htbp]
	\caption{两次建模缺陷的修正与传导（方法论记录）}
	\label{tab:changes}
	\centering
	\zihao{-5}
	\begin{tabular}{p{2.4cm}p{4.6cm}p{4.4cm}p{2.9cm}}
		\toprule
		缺陷 & 根因 & 修正 & 结论变化 \\
		\midrule
		“首批时限物理不可行”
		& \textbf{机队坍缩}（只评估同构机队，把 8 架异构机队压成 2 架 C 型在用）$+$ \textbf{派发退化}（无违规时按能耗打破平局，首轮机位被顺路小架次占用）
		& 候选机队纳入真实混合机队，由真实调度器统一评分；首批专架次按机队槽位轮转播种；派发键加入最小松弛量
		& 准时率由 55.0\% 升至 \textbf{100.0\%}，首批违规 0 箱 \\
		“通信覆盖 100\%”
		& \textbf{服务区间被静默截短} $+$ \textbf{排班未纳入服务窗口} $+$ \textbf{校验器从未执行该项检查}
		& 评估函数显式失败；排班把窗口作为硬约束；一律如实上报中断区间；校验器区分软/硬违规
		& 时间轴覆盖率 $0/20\to$ \textbf{11/20}，并报告 4 架中继缺口 \\		\bottomrule
	\end{tabular}
\end{table}

\textbf{传导链}：问题二修正后，问题三的运输方案与需保障架次全部变化（25 架次中 20 个需保障），问题四的原子单元数由旧方案的 1 个变为 \textbf{3 个}，$K$ 的取值范围随之改变。这说明\textbf{“原子单元数”是下游的导出量，绝不可硬编码}——本文把它作为由程序计算并打印的中间结果，每次上游变化后强制复核。
